%==============================================================================
% Extended Results (analysis)
%
% This file is \input AFTER the existing \section{Extended Results}
% (\label{app:extended}) stub in main.tex. It CONTINUES that section: it opens NO
% new \section and does NOT redefine \label{app:extended}. It gathers the
% per-system failure autopsies, the gbrain "why it leads" mechanism
% (\label{sec:exp:gbrain}, self-defined only; carried here for robustness), the
% gbrain medium failure-mode shift, the two full readiness-gating case write-ups,
% the per-run dollar figures, the scaling edge-count micro-detail, and the judge
% token-budget/temperature/output-cap detail --- all relocated verbatim (lightly
% reflowed) from the condensed body Section~5 (Experiments).
%==============================================================================

\subsection{Judge and cost detail}
\label{app:extended:judge-cost}

\paragraph{Judge configuration.}
Scoring is performed by an LLM judge---Claude Sonnet~4.6 with extended thinking
(a 4000-token reasoning budget, temperature~1.0, up to 8000 output tokens)---
prompted with the question, the gold answer, the evidence artifact IDs, and the
per-sub-criterion rubric, and instructed to award each sub-point independently
before summing. Extended thinking is used deliberately: the judge must itself
perform the multi-hop check the question demands (did the answer cite the
\emph{current} decision rather than the superseded one; is the asserted date
consistent with the bi-temporal record) before it can grade faithfully. The same
judge configuration scores every system, so judge variance is shared across the
comparison.

\paragraph{Per-run cost figures.}
Measured end-to-end cost (judge plus answerer) is roughly \$0.02--\$0.04 per
question and is dominated by those two components for systems with no metered
extraction. On an 11-question small-tier run the judge costs $\approx$\$0.18--\$0.24
and the shared neutral answerer $\approx$\$0.14--\$0.25 (the answerer cost falls on
\texttt{zep-cloud}, \texttt{mem0-platform}, and \texttt{graphify-oss}; \texttt{gbrain}
uses its native \texttt{think} synthesizer and incurs \$0 of shared-answerer cost).
Per-system extraction spend is only metered where it runs on tokens we pay for:
\texttt{graphify-oss} incurs $\approx$\$0.79 of Sonnet extraction at the small tier
and $\approx$\$6.44 at the medium tier---where extraction, not the judge, becomes
the largest single cost component---while \texttt{gbrain} surfaces \$0 at ingest
(capture runs through ZeroEntropy and a local pipeline) and the hosted
\texttt{mem0-platform} and \texttt{zep-cloud} extraction costs are server-side and
not exposed to the harness. Because the harness's \texttt{cost\_per\_query\_usd}
field only captures answerer tokens it issues directly, a system whose synthesizer
runs out-of-process (\texttt{gbrain}'s \texttt{think}) reports \$0 there even though
it does real work; the meaningful cost axis is the system's own
extract-plus-retrieval spend, not the answerer total.

\subsection{Per-system failure autopsies}
\label{app:extended:autopsies}

The category profile explains \emph{why each non-\texttt{gbrain} system lands
where it does}, and the failure-mode labels make the diagnosis concrete.

\paragraph{\texttt{graphify-oss} ($0.206$): coarse extraction $\rightarrow$
retrieval miss.}
The dominant failure mode is \texttt{retrieval\_miss}. The root cause is
extraction granularity: on organizational prose the extractor produces roughly
\emph{one node per artifact} (135 nodes for the 443 medium artifacts; one node
per artifact at small), versus the $\approx 5$ nodes per artifact it produces on
short code samples. Because the graph node \emph{is} the whole document, the
keyword-seeded breadth-first search reaches adjacent whole documents but never
the individual facts \emph{within} a document, so the evidence a question needs
is structurally unreachable. The system also carries no embeddings (retrieval is
pure BFS/DFS over the graph) and no bi-temporal model at all---\texttt{as\_of} is
silently ignored---so C3~($0.028$) and C6~($0.033$) are near zero by
construction. Its only relative strengths, supersession ($0.362$) and
justification chain ($0.406$), come from coarse document-community proximity
rather than from any temporal or relational reasoning.

\paragraph{\texttt{mem0-platform} ($0.221$): a weak undisclosed extractor.}
Here too \texttt{retrieval\_miss} dominates, and the shipped product
extracts with an undisclosed server-side model that we cannot configure. The
$+0.258$ gap to the Sonnet-extractor probe (\S\ref{sec:exp:main}) is the key
evidence that the \emph{extractor}, not the wire protocol or retrieval logic, is
the bottleneck: holding everything else fixed and substituting Sonnet~4.6 more
than doubles the score. The hosted system is competent only on simple decision
provenance (C2~$0.400$) and, interestingly, posts the field's best small-tier C3
($0.317$)---though on 11 questions this is a thin signal. Its $0.000$ on
supersession (C1) and audit replay (C4) shows it does not track facts changing
over time.

\paragraph{\texttt{zep-cloud} ($0.252$): graph completeness helps recall, hurts
multi-hop.}
\texttt{zep-cloud} is the one non-\texttt{gbrain} system that builds a genuine
temporal graph and therefore \emph{engages} the temporal categories rather than
abstaining. It is the strongest of the bottom three on simple recall
(supersession $0.525$, decision provenance $0.625$) but the weakest where the
benchmark's point lives: it scores $0.000$ on audit replay (C4) and only $0.117$
on justification chains (C5). The readiness experiment
(\S\ref{sec:exp:readiness}) exposes the mechanism directly: letting the graph
finish building \emph{raised} simple-recall categories (supersession $+0.53$,
decision provenance $+0.50$) while \emph{lowering} multi-hop ones (contradiction
$-0.33$, justification chain $-0.27$, bi-temporal $-0.23$/$-0.17$). A fuller
graph improves point lookups but dilutes the relevant subgraph at scoring time,
hurting reasoning that must isolate a small set of related edges. A single tier
mean hides this shape entirely---a methodological point we return to in
\S\ref{sec:exp:readiness}.

\subsection{Why \texttt{gbrain} Leads the Public Field}
\label{sec:exp:gbrain}

\texttt{gbrain} leads because of \emph{how it retrieves and synthesizes}, not
because it solves the task. Its retrieval is hybrid---BM25 plus dense vector
plus graph traversal---so it is far less prone to the single-modality blind spots
that sink the others (\texttt{graphify-oss} has no vectors;
\texttt{mem0-platform} has no graph structure exposed to the answerer). Its
native \texttt{think} synthesizer then produces cited, multi-hop answers over the
retrieved set, which is why it wins simple-recall provenance (C2~$0.762$) and
supersession (C1~$0.675$) outright on the small tier. Its small-tier ceiling is
set almost entirely by \emph{retrieval}: the failure modes are
\texttt{retrieval\_miss} with only two hallucinations, meaning that when
it fails on the small tier it is because the evidence never surfaced, not because
it reasoned badly over what it had.

But ``leads'' is not ``solves.'' \texttt{gbrain} scores $0.000$ on small-tier
contradiction (C6) and $0.122$ on bi-temporal (C3): even the strongest public
system cannot reliably detect that two artifacts conflict, or answer what was
true as of a past date. The ``remaining headroom, not the front-runner'' framing
of this gap is the paper's headline result, stated canonically in Main Results
(\S\ref{sec:exp:main}); here we keep only the mechanism behind the ceiling.

\paragraph{Medium-tier failure-mode shift.}
% {\sloppy} absorbs the long unbreakable \texttt{} failure-mode identifiers
% (retrieval_miss / hallucination / wrong_synthesis) so these lines do not run
% over the edge; complements the orchestrator's global \emergencystretch.
{\sloppy
For \texttt{gbrain}, the per-category medium profile is recognizable---provenance
and supersession on top, justification chains at the floor---but with two scale
effects. First, C5 \emph{degrades} ($0.411 \rightarrow 0.216$): longer
justification chains over a larger corpus are markedly harder. Second, the
\emph{failure mode shifts}: at small scale \texttt{gbrain}'s dominant failure is
\texttt{retrieval\_miss}, but at medium scale it becomes \texttt{hallucination}
(ahead of \texttt{retrieval\_miss}), with \texttt{wrong\_synthesis} a distant
third---at organizational scale the system more often surfaces \emph{some} context
and then confidently synthesizes an unsupported answer, rather than surfacing
nothing.\par}

\subsection{Scaling edge-count detail}
\label{app:extended:scaling}

The per-system degradation profiles behind the small-to-medium scaling result
(\S\ref{sec:exp:scaling}) are as follows. \texttt{gbrain} is \emph{scale-robust}:
$0.394 \rightarrow 0.428$ (it even edges up), its hybrid retrieval keeping pace
with the larger corpus. \texttt{mem0-platform} \emph{improves} with scale:
$0.221 \rightarrow 0.303$; the larger corpus yields more extracted memories,
lifting its provenance and contradiction categories (C2~$0.49$, C6~$0.43$).
\texttt{zep-cloud} \emph{degrades}: $0.252 \rightarrow 0.213$; its temporal graph
stays sparse relative to the corpus (it settled at $1000$ edges for $443$
artifacts vs.\ $654$ for $121$), so recall thins out. \texttt{graphify-oss}
\emph{degrades most in relative terms}: $0.206 \rightarrow 0.142$; its
one-node-per-document extraction leaves the within-document facts unreachable,
and the gap widens as documents multiply.

\subsection{Readiness-gating case write-ups}
\label{app:extended:readiness}

The methodological readiness-gating result (\S\ref{sec:exp:readiness}) rests on
two cases where the drain design changed the result.

\paragraph{\texttt{mem0-platform}: $0.141 \rightarrow 0.221$ (the naive drain
understated the system).}
A naive drain that polled \texttt{get\_all} counts returned on a $0\rightarrow
0\rightarrow 0$ plateau---asynchronous extraction had not yet \emph{started}, so
the store looked stably empty---and the benchmark then queried an effectively
empty memory. Replacing it with a probe-\emph{search} drain (requiring four
consecutive \emph{non-zero} stable counts above a 60-second floor) revealed that
extraction actually takes $\approx$420\,s to settle at $\approx$81 memories.
Under the corrected gate the score rose from $0.141$ to $0.221$. Notably the
broken run's slightly higher score on one question came from the answerer
hallucinating a correct-looking answer from the question text with zero
retrieved characters---an artifact of querying an empty store, not real
performance. The lesson is counter-intuitive and important for fair evaluation:
\emph{the naive drain understated the hosted system}.

\paragraph{\texttt{zep-cloud}: a half-built graph times out (the mean moves only
modestly; the shape changes a lot).}
A 600\,s drain ceiling timed out on a half-built graph ($\approx$300 edges); a
proper 1800\,s settle reached $654{+}$ edges. The tier mean moved only modestly
($0.220 \rightarrow 0.252$), but the per-question scores moved \emph{massively}:
supersession $+0.53$ and decision provenance $+0.50$, against contradiction
$-0.33$, justification chain $-0.27$, and bi-temporal $-0.23$/$-0.17$. A more
complete graph helped simple recall and \emph{hurt} multi-hop reasoning (the
larger graph dilutes the relevant subgraph at scoring time)---so a single mean
hides a real change in system behavior, and reporting only the mean would have
masked a genuine completeness-vs-reasoning tradeoff. Reaching the settled state
also required honoring the API's \texttt{Retry-After} (the free plan rate-limits
aggressively) and polling \texttt{graph.get} for propagation to defeat a
delete$\rightarrow$create$\rightarrow$add 404 race during store reset.
